\begin{table}[h!]
\centering
\caption{Ablation Study of Proposed Method on Solar Filament Segmentation Benchmark}
\label{tab:ablation}
\begin{tabular}{lcccc}
\hline
\textbf{Method Variant} & \textbf{PQ $\uparrow$} & \textbf{SQ $\uparrow$} & \textbf{RQ $\uparrow$} & \textbf{Mean Dice $\uparrow$} \\
\hline
1. Baseline (Threshold 0.50 only) & 0.3576 & 0.6811 & 0.5251 & 0.6764 \\
2. + Calibrated Threshold (0.52) & 0.3559 & 0.6805 & 0.5231 & 0.6760 \\
3. + Min-Area Noise Filter (260px) & 0.4127 & 0.6812 & 0.6059 & 0.6771 \\
4. + Curvilinear Grouping (r=10) & 0.4251 & 0.6769 & 0.6280 & 0.6781 \\
5. + Dual-Backbone Ensemble and Zoom TTA (v7 Peak) & 0.4231 & 0.6735 & 0.6282 & 0.6756 \\
6. + Tri-Backbone Diversity Ensemble (v9 Peak) & 0.4199 & 0.6727 & 0.6242 & 0.6741 \\
7. + Thickness-Aware Spine Recovery (v12 Peak) & 0.4230 & 0.6745 & 0.6271 & 0.6761 \\
8. + Full Network Station Calibration and TAC-SR v2 (v20 Peak) & 0.4318 & 0.6730 & 0.6416 & 0.6755 \\
9. + Extended Grid and TAC-SR (p25) (v21 Baseline) & 0.4302 & 0.6731 & 0.6391 & 0.6749 \\
10. + Full 67-Obs PQ Checkpointing and 5.5 fil/img Calibration (v22 SOTA) & 0.4346 & 0.6732 & 0.6455 & 0.6741 \\
\hline
\end{tabular}
\end{table}
